\subsection{Complements on Characteristic Polynomials}

Let L be a commutative ring and M a free L-module of finite rank m. If u is an endomorphism of M and r a natural number, then we denote by \(c_{r}(u)\) the trace of the endomorphism \(\wedge^{r}(u)\) of the free L-module \(\wedge^{r}(\mathrm{M})\) . In particular, we have

\[
c _ {0} (u) = 1, \quad c _ {1} (u) = \mathrm{Tr} (u), \quad c _ {m} (u) = \det (u),\tag{1}
\]

and \(c_{r}(u) = 0\) for r \textgreater{} m. By Proposition 7 of III, §8, No.~4, p.~527, the mapping \(u \mapsto \det(u)\) from End(M) to L is a homogeneous polynomial mapping of degree m (IV, §5, No.~9, p.~55). More generally, for every integer r such that \(0 \leqslant r \leqslant m\) , the mapping \(c_{r}\) from End(M) to L is a homogeneous polynomial mapping of degree r; this follows from Proposition 10 of III, §8, No.~5, p.~529.

Let \(u\) be an endomorphism of \(M\) and \(\overline{u}\) the endomorphism of the \(L[X]\) -module \(M[X] = M \otimes_{L} L[X]\) deduced from \(u\) by extension of scalars (II, §5, No.~1, p.~277). Recall (III, §8, No.~11, p.~541, Definition 3 and (50)) that the characteristic polynomial of \(u\) is the determinant \(\chi_u(X)\) of the \(L[X]\) -endomorphism \(X - \overline{u}\) of \(M[X]\) and that we have the relation

\[
\chi_ {u} (\mathrm{X}) = \sum_ {r = 0} ^ {m} (- 1) ^ {r} c _ {r} (u) \mathrm{X} ^ {m - r}.\tag{2}
\]

PROPOSITION 1. --- Let L be a commutative ring, M a free L-module of finite rank \(m \geqslant 1\) , and u an endomorphism of M. There exists a unique endomorphism \(\widetilde{u}\) of M satisfying the relation

\[
\widetilde {u} (x) \wedge w = x \wedge \wedge^ {m - 1} (u) (w)\tag{3}
\]

for \(x\in \mathrm{M}\) and \(w\in \bigwedge^{m - 1}(\mathrm{M})\) . Moreover, we have the relations

\[
u \circ \widetilde {u} = \widetilde {u} \circ u = \det (u) _ {\mathrm{M}},\tag{4}
\]

\[
\det (\widetilde {u}) = \det (u) ^ {m - 1},\tag{5}
\]

\[
\widetilde {u} = \sum_ {r = 0} ^ {m - 1} (- 1) ^ {r} c _ {m - 1 - r} (u) u ^ {r}.\tag{6}
\]

Lemma 1. --- Let p be an integer such that \(0 \leqslant p \leqslant m\) . For any w in \(\wedge^{p}(\mathrm{M})\) , let \(h_{p}(w)\) be the linear mapping \(w' \mapsto w \wedge w'\) from \(\wedge^{m-p}(\mathrm{M})\) to \(\wedge^{m}(\mathrm{M})\) . The linear mapping \(h_{p}: w \mapsto h_{p}(w)\) from \(\wedge^{p}(\mathrm{M})\) to \(\operatorname{Hom}_{\mathrm{L}}(\wedge^{m-p}(\mathrm{M}), \wedge^{m}(\mathrm{M}))\) is an isomorphism.

Let \((e_{i})_{i\in\mathrm{I}}\) be a basis of M; we endow the set I with a total order. For any subset J of I, set \(e_{\mathrm{J}}=e_{i_{1}}\wedge\cdots\wedge e_{i_{r}}\) , where \((i_{1},\ldots,i_{r})\) is the sequence of elements of J in increasing order. The L-module \(\wedge^{m-p}(\mathrm{M})\) admits as a basis the elements \(e_{\mathrm{S}}\) , where S runs through the set of subsets of I with m-p elements; \(\wedge^{m}(\mathrm{M})\) has \(\{e_{\mathrm{I}}\}\) as a basis. Consequently, there exists a basis of \(\mathrm{Hom}_{\mathrm{L}}(\wedge^{m-p}(\mathrm{M}),\wedge^{m}(\mathrm{M}))\) consisting of linear mappings \(e_{\mathrm{J}}^{*}\) characterized by the formula

\[
e _ {\mathrm{J}} ^ {*} (e _ {\mathrm{S}}) = \left\{ \begin{array}{l l} e _ {\mathrm{I}} & \text {if } I = J\cup S, \\ 0 & \text {otherwise}, \end{array} \right.\tag{7}
\]

where J runs through the set of subsets of I with p elements. It follows from formula (20) of III, §7, No.~8, p.~519 that for every subset J of I with p elements, we have \(h_{p}(e_{\mathrm{J}}) \in \{e_{\mathrm{J}}^{*}, -e_{\mathrm{J}}^{*}\}\) ; since the elements \(e_{J}\) form a basis of \(\Lambda^{p}(\mathrm{M})\) , the linear mapping \(h_{p}\) is bijective.

Let us now prove Proposition 1. Let \(u\) and \(\widetilde{u}\) be endomorphisms of M. Relation (3) is equivalent to

\[
h _ {1} \circ \widetilde {u} = \operatorname{Hom} (\wedge^ {m - 1} (u) \cdot 1 _ {\bigwedge^ {m} (\mathrm{M})}) \circ h _ {1};\tag{8}
\]

the mapping \(h_{1}\) is an isomorphism from M to \(\operatorname{Hom}_{\mathrm{L}}(\wedge^{m-1}(\mathrm{M}),\wedge^{m}(\mathrm{M}))\) by Lemma 1. Consequently, for every endomorphism u of M, there exists a unique endomorphism \(\widetilde{u}\) of M satisfying relation (3).

Let \(x_{1},\ldots,x_{m}\) be elements of M. Let us replace x with \(u(x_{1})\) and w with \(x_{2}\wedge\cdots\wedge x_{m}\) in (3); we obtain

\[
\widetilde {u} (u (x _ {1})) \wedge x _ {2} \wedge \dots \wedge x _ {m} = u (x _ {1}) \wedge \dots \wedge u (x _ {m}) = \det (u) x _ {1} \wedge \dots \wedge x _ {m}.
\]

Consequently, \(h_1(\widetilde{u} \circ u(x_1)) = h_1(\det(u)x_1)\) , which gives the relation \(\widetilde{u} \circ u = \det(u)_\mathrm{M}\) by Lemma 1.

We denote by U the endomorphism \(X - \overline{u}\) of the L{[}X{]}-module M{[}X{]} (VIII, p.~335). By the above applied to U, there exists an endomorphism \(\widetilde{U}\) of the L{[}X{]}-module M{[}X{]} that satisfies the relations

\[
\widetilde {\mathrm{U}} (x _ {1}) \wedge x _ {2} \wedge \dots \wedge x _ {m} = x _ {1} \wedge (\mathrm{X} x _ {2} - u (x _ {2})) \wedge \dots \wedge (\mathrm{X} x _ {m} - u (x _ {m}))\tag{9}
\]

for \(x_{1},\ldots ,x_{m}\) in M and

\[
\widetilde {\mathrm{U}} \circ \mathrm{U} = \det (\mathrm{X} - \overline {{u}}) _ {\mathrm{M} [ \mathrm{X} ]}.\tag{10}
\]

Let us view \(\widetilde{\mathbf{U}}\) as an element of \(\operatorname{End}(\mathbf{M})[\mathbf{X}]\) (VIII, p.~9); by formula (9) and Lemma 1, it has degree \(\leqslant m - 1\) , so we can write it as

\[
\widetilde {\mathrm{U}} = \sum_ {r = 0} ^ {m - 1} (- 1) ^ {r} u _ {r} X ^ {m - 1 - r},\tag{11}
\]

where the \(u_{r}\) are endomorphisms of M. By formula (2), relation (10) gives the equality

\[
\left(\sum_ {r = 0} ^ {m - 1} (- 1) ^ {r} u _ {r} \mathrm{X} ^ {m - 1 - r}\right) (\mathrm{X} - u) = \sum_ {r = 0} ^ {m} (- 1) ^ {r} c _ {r} (u) \mathrm{X} ^ {m - r}\tag{12}
\]

in the ring \(\operatorname{End}(M)[X]\) . By identifying the coefficients of the monomials in X on each side, we obtain the relations

\[
u _ {r} + u _ {r - 1} \circ u = c _ {r} (u) _ {\mathrm{M}} \qquad \mathrm{for} 1 \leqslant r \leqslant m - 1\tag{13}
\]

and

\[
u _ {0} = c _ {0} (u) _ {\mathrm{M}}, \quad u _ {m - 1} \circ u = c _ {m} (u) _ {\mathrm{M}}.\tag{14}
\]

From this, we deduce

\[
u _ {m - 1} = \sum_ {r = 0} ^ {m - 1} (- 1) ^ {r} c _ {m - 1 - r} (u) u ^ {r}.\tag{15}
\]

Now, when we identify the constant terms, equalities (9) and (11) imply \(u_{m-1} = \widetilde{u}\) ; formula (6) follows.

In particular, \(\widetilde{u}\) belongs to the subalgebra of \(\operatorname{End}(\mathbf{M})\) generated by \(u\) and therefore commutes with \(u\) . We have already established the relation \(\widetilde{u} \circ u = \det(u)_{\mathrm{M}}\) ; formula (4) follows.

Finally, let \(x_{1},\ldots,x_{m}\) be elements of M. We replace x with \(x_{1}\) and w with \(\widetilde{u}(x_{2})\wedge\cdots\wedge\widetilde{u}(x_{m})\) in formula (3). This gives

\[
\begin{array}{c} \widetilde {u} (x _ {1}) \wedge \widetilde {u} (x _ {2}) \wedge \dots \wedge \widetilde {u} (x _ {m}) = x _ {1} \wedge u \circ (\widetilde {u} (x _ {2})) \wedge \dots \wedge u \circ (\widetilde {u} (x _ {m})) \\ = \det (u) ^ {m - 1} x _ {1} \wedge x _ {2} \wedge \dots \wedge x _ {m} \end{array}
\]

and therefore formula (5).

Remarks. --- 1) The endomorphism \(\widetilde{u}\) of M coincides with what we called the cotranspose of \(u\) in III, §8, No.~6, p.~532.

\begin{enumerate}
\def\labelenumi{\arabic{enumi})}
\setcounter{enumi}{1}
\item
  From formulas (1), (2), (4), and (6), we deduce the relation \(\chi_u(u) = 0\) and therefore another proof of the Cayley-Hamilton theorem (III, §8, No.~11, p.~541).
\item
  Since the mapping \(c_{r}\) from End(M) to L is a homogeneous polynomial mapping of degree r for \(0 \leqslant r \leqslant m - 1\) , it follows from formula (6) that the mapping \(u \mapsto \widetilde{u}\) from End(M) to End(M) is a homogeneous polynomial mapping of degree m - 1.
\end{enumerate}

Let B be an algebra over the ring L; suppose that B is a free L-module of rank \(m \geqslant 1\) , and identify L with the subring \(L \cdot 1\) of B. Let b be an element of B. We apply the above to the endomorphism \(\gamma(b) : x \mapsto bx\) of the L-module B. Set \(\gamma_{r}(b) = c_{r}(\gamma(b))\) for \(0 \leqslant r \leqslant m\) ; we have, in particular, \(\gamma_{m}(b) = \mathrm{N}_{\mathrm{B/L}}(b)\) (III, §9, No.~3, p.~543). The characteristic polynomial of b (loc. cit.) can be written as

\[
\mathrm{Pc} _ {\mathrm{B/L}} (b; \mathrm{X}) = \sum_ {r = 0} ^ {m} (- 1) ^ {r} \gamma_ {r} (b) \mathrm{X} ^ {m - r}.\tag{16}
\]

Since the mapping \(\gamma\) from B to \(\operatorname{End}_{\mathrm{L}}(\mathrm{B})\) is L-linear, the mapping \(\gamma_{r}\) from B to L is a homogeneous polynomial mapping of degree r. In particular, the mapping \(b \mapsto \mathrm{N}_{\mathrm{B}/\mathrm{L}}(b)\) from B to L is a homogeneous polynomial mapping of degree m.

For any element \(b\) of B, set

\[
\widetilde {b} = \sum_ {r = 0} ^ {m - 1} (- 1) ^ {r} \gamma_ {m - 1 - r} (b) b ^ {r}.\tag{17}
\]

By Proposition 1, the linear mapping \(\gamma (\widetilde{b})\) from B to B is the cotranspose of the mapping \(\gamma (b)\) ; from this, we deduce

\[
b \widetilde {b} = \widetilde {b} b = \mathrm{N} _ {\mathrm{B/L}} (b).\tag{18}
\]

Moreover, the mapping \(b \mapsto \widetilde{b}\) from B to B is a homogeneous polynomial mapping of degree \(m - 1\) .

